\documentclass[10pt,a4paper]{amsart}
\usepackage[margin=19mm]{geometry}
\usepackage{amsmath,amssymb,mathtools}
\usepackage[scr=rsfs,cal=euler]{mathalfa}
\usepackage{xfrac}
\usepackage[unicode,hidelinks,bookmarks=false]{hyperref}
\newcommand{\ie}{i.e.\ }
\newcommand{\cf}{cf.\ }
\newcommand{\resp}{respectively\ }
\newcommand{\Subsection}{\subsection}
\providecommand{\nobreakdash}{\nobreak\hbox{-}}
\setlength{\emergencystretch}{2em}
\multlinegap0pt
\usepackage{amssymb}
\usepackage{bm}
\usepackage[shortlabels]{enumitem}
\newtheorem{lemma}{Lemma}
\newcommand{\mc}{\mathcal}
\title{Counting $2\times 2$ matrices with fixed determinant and bounded coefficients}
\author{Kavita Dhanda, Alan Haynes, Silmi Prasala}
\date{Source: arXiv:2509.16890v1 (CC BY 4.0)}
\begin{document}
\maketitle

\noindent\emph{Source and licence.} Counting $2\times 2$ matrices with fixed determinant and bounded coefficients. Version arXiv:2509.16890v1 (CC BY 4.0), distributed by arXiv under the Creative Commons Attribution 4.0 International licence, http://creativecommons.org/licenses/by/4.0/, as recorded in the arXiv licence metadata for this exact version and shown on the abstract page https://arxiv.org/abs/2509.16890v1. Full licence notice: \texttt{licenses/arxiv-2509.16890-CC-BY-4.0.txt}.\par\smallskip

\noindent\textit{Editorial scope.} Counted unit: the article's lemma lem.p-adic1 (source0.tex lines 208--211). The article's complete original proof of it is delivered with it (source0.tex lines 212--220, one \texttt{proof} environment of 9 lines). No mathematics is written, repaired or re-derived: the statement and the proof are the article's own lines, in the article's own notation, and no retained line is substituted or re-flowed.  No further source-local context is needed: every cross-reference of every retained line resolves inside the two delivered spans.  Nothing was omitted from any retained span, so the package carries no omission record.  No retained line contains a citation, so the wrapper carries no bibliography and no bibliography entry.  No retained line contains a figure environment, an image-inclusion command or drawing code, so no image is delivered and the package declares no asset dependency.  Editorial formatting only, in addition: an AMS \texttt{amsart} preamble that restates the article's own declarations and macros used by the retained lines, namely the environments \texttt{lemma} on separate counters, together with the standard packages those lines need.  The retained environments print in this document as Lemma 1 in order of appearance, whereas the article prints the same items with the numbering of its own full document.  The complete native source of the article as distributed in the arXiv e-print for this exact version is bundled unchanged as \texttt{sources/185419/source0.tex}.
\begin{lemma}\label{lem.p-adic1}
Let $p$ be prime and suppose that $n=p$ or that $n=p^2$. For every integer $c$ satisfying $0\le c< n$, there exist integers $a$ and $q$ satisfying $|a|\le n^{1/2}$, $0<q< n^{1/2}$, and
\[aq^{-1}=c ~\mathrm{mod}~ n.\]
\end{lemma}

\begin{proof}[Proof of lemma]
Let
\[\mc{A}=\left\{(b,r): 0\le b\le n^{1/2}, 0\le r<n^{1/2}\right\}.\]
It is easy to check that, whether $n=p$ or $p^2$,
\[|\mc{A}|\ge n+1.\]
Therefore, by the Dirichlet pigeonhole principle, there are distinct $(a_1,q_1),(a_2,q_2)\in\mc{A}$ with
\[a_1-q_1c=a_2-q_2c~\mathrm{mod}~n.\]
Take $a=a_1-a_2$ and $q=q_1-q_2$ and assume without loss of generality, by switching the signs if necessary, that $q\ge 0$. If $q=0~\mathrm{mod}~p$ then, since $q<n^{1/2}\le p$, it must be the case that $q=0$. Since $|a|\le n^{1/2}$, this forces $a=0$ as well, which is a contradiction. Therefore $q>0$, and it is invertible modulo $n$, which completes the proof of the lemma.
\end{proof}

\end{document}
